\begin{frame}{왜 $\sigma$ 도 상태 함수가 아닌가}
  ``$\sigma$ 도 두 번째 헤드에서 $\log\sigma(s)$ 로 출력하면 안
  되나?'' --- 가능하지만, 표준 PPO 구현은 \textbf{전역 버전}이 일반적:
  \begin{itemize}
    \item \textbf{(1) 파라미터 경제성}: 상태별 $\sigma$ 는 64$\to$2 출력
      헤드가 하나 더 늘지만 이득 제한적 --- ``얼마나 대담하게
      탐색할지''는 대부분 \emph{상태의} 성질보다 \emph{과업 전체}의
      성질.
    \item \textbf{(2) 안정성}: $\log\sigma$ 는 엔트로피 보상에 직접
      들어가는 값(11.2절의 \texttt{-0.01 * d.entropy()}). 상태마다
      $\sigma$ 가 들썩이면 엔트로피 항도 들썩여 \emph{그래디언트를
      불안정}하게 만듦.
  \end{itemize}
  \vskip0.6em
  {\footnotesize ``간단한 것부터, 안 되면서 복잡하게''라는 11.1절
  클리핑 하이퍼파라미터 선택의 교훈이 또 한 번 등장.}
\end{frame}
